% Appendices of the arXiv version. Tables come from tables/ (make_tables.py);
% App. C and App. D's text are ported from the camera-ready.

\section{Per-seed results}
\label{app:perseed}

Table~\ref{tab:perseed} lists every run behind Table~\ref{tab:results}: all 36
arm\,$\times$\,seed\,$\times$\,domain evaluations, with the hardware each run was trained on.
Seed 42 of the frozen un-augmented arm was trained end to end, and seeds 43 and 44 retrain its
head on cached features. With the backbone frozen and the inputs un-augmented, the head sees
the same features either way, and the published seed-42 head reproduces its errors on the
cache to within $0.0005$ (App.~\ref{app:revision}). Table~\ref{tab:frozen} instead retrains the
head for all three seeds, so its row for the same configuration differs slightly from this
arm's mean, and the body quotes Table~\ref{tab:frozen} because its comparisons are made there.
Each run
keeps the checkpoint with the best \emph{synthetic} validation rotation error, chosen without
reference to either real domain.

\begin{table}[ht]
  \caption{Per-seed results behind Table~\ref{tab:results}. $\Etrans$ range-normalised, $\Erot$ in rad (chance level 2.207). \emph{Head only}: the head retrained on cached frozen features.}
  \label{tab:perseed}
  \centering
  \footnotesize
  \begin{tabular}{ll c l l cc}
    \toprule
    Backbone & Aug. & Seed & Trained on & Domain & $\Etrans$ & $\Erot$ \\
    \midrule
    frozen & -- & 42 & Colab & synthetic (val) & 0.066 & 1.053 \\
    frozen & -- & 42 & Colab & lightbox & 0.151 & 1.752 \\
    frozen & -- & 42 & Colab & sunlamp & 0.253 & 2.048 \\
    frozen & -- & 43 & head only & synthetic (val) & 0.064 & 1.045 \\
    frozen & -- & 43 & head only & lightbox & 0.150 & 1.746 \\
    frozen & -- & 43 & head only & sunlamp & 0.263 & 1.952 \\
    frozen & -- & 44 & head only & synthetic (val) & 0.066 & 1.036 \\
    frozen & -- & 44 & head only & lightbox & 0.148 & 1.791 \\
    frozen & -- & 44 & head only & sunlamp & 0.246 & 2.013 \\
    \midrule
    frozen & \checkmark & 42 & Colab & synthetic (val) & 0.071 & 1.180 \\
    frozen & \checkmark & 42 & Colab & lightbox & 0.118 & 1.718 \\
    frozen & \checkmark & 42 & Colab & sunlamp & 0.306 & 1.864 \\
    frozen & \checkmark & 43 & M5 & synthetic (val) & 0.073 & 1.205 \\
    frozen & \checkmark & 43 & M5 & lightbox & 0.126 & 1.734 \\
    frozen & \checkmark & 43 & M5 & sunlamp & 0.328 & 1.883 \\
    frozen & \checkmark & 44 & M5 & synthetic (val) & 0.074 & 1.169 \\
    frozen & \checkmark & 44 & M5 & lightbox & 0.118 & 1.673 \\
    frozen & \checkmark & 44 & M5 & sunlamp & 0.279 & 1.881 \\
    \midrule
    layer 4 & -- & 42 & Colab & synthetic (val) & 0.039 & 0.306 \\
    layer 4 & -- & 42 & Colab & lightbox & 0.200 & 1.169 \\
    layer 4 & -- & 42 & Colab & sunlamp & 0.299 & 1.476 \\
    layer 4 & -- & 43 & Colab & synthetic (val) & 0.035 & 0.326 \\
    layer 4 & -- & 43 & Colab & lightbox & 0.152 & 1.182 \\
    layer 4 & -- & 43 & Colab & sunlamp & 0.352 & 1.619 \\
    layer 4 & -- & 44 & Colab & synthetic (val) & 0.035 & 0.318 \\
    layer 4 & -- & 44 & Colab & lightbox & 0.154 & 1.147 \\
    layer 4 & -- & 44 & Colab & sunlamp & 0.283 & 1.506 \\
    \midrule
    layer 4 & \checkmark & 42 & Colab & synthetic (val) & 0.041 & 0.352 \\
    layer 4 & \checkmark & 42 & Colab & lightbox & 0.152 & 0.986 \\
    layer 4 & \checkmark & 42 & Colab & sunlamp & 0.405 & 1.592 \\
    layer 4 & \checkmark & 43 & Colab & synthetic (val) & 0.044 & 0.302 \\
    layer 4 & \checkmark & 43 & Colab & lightbox & 0.165 & 0.997 \\
    layer 4 & \checkmark & 43 & Colab & sunlamp & 0.399 & 1.354 \\
    layer 4 & \checkmark & 44 & Colab & synthetic (val) & 0.042 & 0.311 \\
    layer 4 & \checkmark & 44 & Colab & lightbox & 0.132 & 0.977 \\
    layer 4 & \checkmark & 44 & Colab & sunlamp & 0.261 & 1.459 \\
    \bottomrule
  \end{tabular}
\end{table}


\section{Conformal coverage, all arms and domains}
\label{app:conformal}

Split-conformal regions are calibrated per arm on 1{,}000 held-out synthetic images, used
neither for training nor for checkpoint selection, at nominal 90\% coverage with the
finite-sample-corrected quantile. Coverage carries Clopper--Pearson 95\%
intervals~\citep{clopper1934use}. $\hat{q}$ is the calibrated region size: a keep-out sphere
radius for translation, a cone half-angle for rotation. \emph{Inflation} is the factor by which
$\hat{q}$ would have to grow to restore 90\% coverage on a domain, so it measures directly how
wrong the calibration is. It sits near $1$ in-domain and reaches $1.2$--$13\times$ on the real
domains, worst on the sharper layer-4 arms. A low factor on the frozen arms is not good news:
their cones are already too wide in-domain to be informative. There are two kinds of
in-domain row: \emph{in-domain test} rows are a further 1{,}000 fresh synthetic images, and
\emph{Synthetic (val)} rows are the 2{,}000-image validation split used for checkpoint
selection, kept as a second in-domain control that is not fully independent.\footnote{The camera-ready labelled
these two ``synthetic (val)'' and ``synthetic (selection split)''.}

Table~\ref{tab:conformal-abs} uses absolute scores: Euclidean metres for translation, geodesic
radians for rotation. Table~\ref{tab:conformal-rel} repeats the translation analysis under the
range-normalised score. Table~\ref{tab:conformal-revision} gives rotation coverage for the
ablation arms and for every frozen-backbone head.

\begin{table}[ht]
  \caption{Conformal coverage at nominal 90\%, absolute scores, for every run in Table~\ref{tab:results}. Coverage in \% with Clopper--Pearson 95\% intervals, $\hat{q}_T$ in metres and $\hat{q}_R$ in radians.}
  \label{tab:conformal-abs}
  \centering
  \scriptsize
  \setlength{\tabcolsep}{3pt}
  \begin{tabular}{ll c l ccc ccc}
    \toprule
    & & & & \multicolumn{3}{c}{Translation (sphere)} & \multicolumn{3}{c}{Rotation (cone)} \\
    \cmidrule(lr){5-7} \cmidrule(lr){8-10}
    Backbone & Aug. & Seed & Domain & Cov. & $\hat{q}_T$ & Infl. & Cov. & $\hat{q}_R$ & Infl. \\
    \midrule
    frozen & -- & 42 & in-domain test & 87.7 (85.5--89.7) & 0.712 & 1.11 & 90.0 (88.0--91.8) & 2.299 & 1.01 \\
    frozen & -- & 42 & synthetic (val) & 88.0 (86.5--89.4) & 0.712 & 1.08 & 90.5 (89.1--91.7) & 2.299 & 0.99 \\
    frozen & -- & 42 & lightbox & 52.1 (50.9--53.3) & 0.712 & 2.48 & 71.0 (69.9--72.1) & 2.299 & 1.24 \\
    frozen & -- & 42 & sunlamp & 29.2 (27.5--30.9) & 0.712 & 4.76 & 58.2 (56.3--60.0) & 2.299 & 1.28 \\
    frozen & -- & 43 & in-domain test & 90.7 (88.7--92.4) & 0.751 & 0.96 & 88.9 (86.8--90.8) & 2.215 & 1.04 \\
    frozen & -- & 43 & synthetic (val) & 90.5 (89.1--91.7) & 0.751 & 0.98 & 89.1 (87.7--90.5) & 2.215 & 1.03 \\
    frozen & -- & 43 & lightbox & 55.4 (54.2--56.6) & 0.751 & 2.32 & 67.0 (65.9--68.1) & 2.215 & 1.30 \\
    frozen & -- & 43 & sunlamp & 32.0 (30.3--33.8) & 0.751 & 4.95 & 59.1 (57.3--61.0) & 2.215 & 1.31 \\
    frozen & -- & 44 & in-domain test & 88.3 (86.1--90.2) & 0.704 & 1.05 & 88.7 (86.6--90.6) & 2.159 & 1.03 \\
    frozen & -- & 44 & synthetic (val) & 88.6 (87.1--90.0) & 0.704 & 1.06 & 88.4 (87.0--89.8) & 2.159 & 1.04 \\
    frozen & -- & 44 & lightbox & 53.4 (52.2--54.6) & 0.704 & 2.58 & 63.6 (62.4--64.7) & 2.159 & 1.33 \\
    frozen & -- & 44 & sunlamp & 30.6 (28.9--32.3) & 0.704 & 4.72 & 53.4 (51.6--55.3) & 2.159 & 1.36 \\
    \midrule
    frozen & \checkmark & 42 & in-domain test & 89.3 (87.2--91.1) & 0.793 & 1.01 & 88.7 (86.6--90.6) & 2.403 & 1.02 \\
    frozen & \checkmark & 42 & synthetic (val) & 89.9 (88.5--91.2) & 0.793 & 1.01 & 89.5 (88.0--90.8) & 2.403 & 1.02 \\
    frozen & \checkmark & 42 & lightbox & 68.3 (67.2--69.4) & 0.793 & 1.84 & 75.1 (74.1--76.2) & 2.403 & 1.19 \\
    frozen & \checkmark & 42 & sunlamp & 26.3 (24.6--27.9) & 0.793 & 5.36 & 71.0 (69.3--72.7) & 2.403 & 1.20 \\
    frozen & \checkmark & 43 & in-domain test & 89.7 (87.6--91.5) & 0.828 & 1.01 & 90.9 (88.9--92.6) & 2.415 & 0.98 \\
    frozen & \checkmark & 43 & synthetic (val) & 89.6 (88.2--91.0) & 0.828 & 1.01 & 88.9 (87.4--90.2) & 2.415 & 1.03 \\
    frozen & \checkmark & 43 & lightbox & 66.9 (65.7--68.0) & 0.828 & 1.91 & 75.5 (74.4--76.5) & 2.415 & 1.18 \\
    frozen & \checkmark & 43 & sunlamp & 23.4 (21.8--25.0) & 0.828 & 5.16 & 70.0 (68.3--71.7) & 2.415 & 1.20 \\
    frozen & \checkmark & 44 & in-domain test & 89.8 (87.8--91.6) & 0.839 & 1.00 & 91.1 (89.2--92.8) & 2.372 & 0.97 \\
    frozen & \checkmark & 44 & synthetic (val) & 89.8 (88.4--91.1) & 0.839 & 1.01 & 89.5 (88.1--90.8) & 2.372 & 1.01 \\
    frozen & \checkmark & 44 & lightbox & 71.1 (70.0--72.2) & 0.839 & 1.72 & 76.3 (75.2--77.3) & 2.372 & 1.18 \\
    frozen & \checkmark & 44 & sunlamp & 30.1 (28.4--31.9) & 0.839 & 4.69 & 69.8 (68.1--71.5) & 2.372 & 1.21 \\
    \midrule
    layer 4 & -- & 42 & in-domain test & 89.3 (87.2--91.1) & 0.410 & 1.04 & 88.8 (86.7--90.7) & 0.460 & 1.06 \\
    layer 4 & -- & 42 & synthetic (val) & 89.6 (88.2--90.9) & 0.410 & 1.01 & 87.2 (85.7--88.7) & 0.460 & 1.10 \\
    layer 4 & -- & 42 & lightbox & 34.1 (32.9--35.2) & 0.410 & 7.92 & 21.5 (20.5--22.5) & 0.460 & 5.52 \\
    layer 4 & -- & 42 & sunlamp & 18.2 (16.8--19.6) & 0.410 & 10.92 & 9.1 (8.1--10.3) & 0.460 & 6.07 \\
    layer 4 & -- & 43 & in-domain test & 90.2 (88.2--92.0) & 0.378 & 0.99 & 89.1 (87.0--91.0) & 0.540 & 1.04 \\
    layer 4 & -- & 43 & synthetic (val) & 89.2 (87.8--90.5) & 0.378 & 1.02 & 89.2 (87.8--90.5) & 0.540 & 1.05 \\
    layer 4 & -- & 43 & lightbox & 35.4 (34.3--36.6) & 0.378 & 6.41 & 26.7 (25.6--27.7) & 0.540 & 4.76 \\
    layer 4 & -- & 43 & sunlamp & 13.8 (12.6--15.2) & 0.378 & 13.42 & 9.0 (8.0--10.1) & 0.540 & 5.23 \\
    layer 4 & -- & 44 & in-domain test & 89.9 (87.9--91.7) & 0.368 & 1.01 & 89.7 (87.6--91.5) & 0.526 & 1.02 \\
    layer 4 & -- & 44 & synthetic (val) & 88.4 (87.0--89.8) & 0.368 & 1.05 & 90.5 (89.1--91.7) & 0.526 & 0.99 \\
    layer 4 & -- & 44 & lightbox & 32.9 (31.8--34.1) & 0.368 & 6.52 & 26.5 (25.4--27.5) & 0.526 & 4.71 \\
    layer 4 & -- & 44 & sunlamp & 18.0 (16.6--19.5) & 0.368 & 11.73 & 12.9 (11.7--14.2) & 0.526 & 5.30 \\
    \midrule
    layer 4 & \checkmark & 42 & in-domain test & 89.8 (87.8--91.6) & 0.425 & 1.01 & 91.2 (89.3--92.9) & 0.576 & 0.97 \\
    layer 4 & \checkmark & 42 & synthetic (val) & 89.8 (88.3--91.0) & 0.425 & 1.01 & 90.0 (88.7--91.3) & 0.576 & 1.00 \\
    layer 4 & \checkmark & 42 & lightbox & 44.3 (43.1--45.5) & 0.425 & 6.45 & 37.2 (36.1--38.4) & 0.576 & 3.82 \\
    layer 4 & \checkmark & 42 & sunlamp & 10.0 (8.9--11.2) & 0.425 & 12.79 & 12.6 (11.4--13.9) & 0.576 & 4.88 \\
    layer 4 & \checkmark & 43 & in-domain test & 90.4 (88.4--92.2) & 0.494 & 0.98 & 91.1 (89.2--92.8) & 0.483 & 0.96 \\
    layer 4 & \checkmark & 43 & synthetic (val) & 91.3 (90.0--92.5) & 0.494 & 0.95 & 89.8 (88.3--91.0) & 0.483 & 1.01 \\
    layer 4 & \checkmark & 43 & lightbox & 47.8 (46.6--49.0) & 0.494 & 6.15 & 32.6 (31.5--33.7) & 0.483 & 4.85 \\
    layer 4 & \checkmark & 43 & sunlamp & 10.7 (9.6--11.9) & 0.494 & 10.80 & 15.3 (14.0--16.7) & 0.483 & 5.59 \\
    layer 4 & \checkmark & 44 & in-domain test & 91.0 (89.1--92.7) & 0.460 & 0.98 & 91.4 (89.5--93.1) & 0.544 & 0.92 \\
    layer 4 & \checkmark & 44 & synthetic (val) & 91.0 (89.6--92.2) & 0.460 & 0.97 & 90.8 (89.4--92.0) & 0.544 & 0.96 \\
    layer 4 & \checkmark & 44 & lightbox & 51.0 (49.8--52.2) & 0.460 & 5.10 & 38.0 (36.9--39.2) & 0.544 & 4.19 \\
    layer 4 & \checkmark & 44 & sunlamp & 24.4 (22.9--26.1) & 0.460 & 9.21 & 15.4 (14.1--16.8) & 0.544 & 5.03 \\
    \bottomrule
  \end{tabular}
\end{table}

\begin{table}[ht]
  \caption{Conformal coverage at nominal 90\%, translation scored with the range-normalised error $\Etrans$.}
  \label{tab:conformal-rel}
  \centering
  \scriptsize
  \setlength{\tabcolsep}{4pt}
  \begin{tabular}{ll c l ccc}
    \toprule
    Backbone & Aug. & Seed & Domain & Coverage & $\hat{q}$ (relative) & Infl. \\
    \midrule
    frozen & -- & 42 & in-domain test & 89.4 (87.3--91.2) & 0.1189 & 1.01 \\
    frozen & -- & 42 & synthetic (val) & 89.3 (87.9--90.7) & 0.1189 & 1.03 \\
    frozen & -- & 42 & lightbox & 49.6 (48.4--50.8) & 0.1189 & 2.42 \\
    frozen & -- & 42 & sunlamp & 21.7 (20.2--23.3) & 0.1189 & 3.85 \\
    frozen & -- & 43 & in-domain test & 90.3 (88.3--92.1) & 0.1217 & 0.99 \\
    frozen & -- & 43 & synthetic (val) & 91.2 (89.9--92.4) & 0.1217 & 0.94 \\
    frozen & -- & 43 & lightbox & 51.7 (50.5--52.9) & 0.1217 & 2.38 \\
    frozen & -- & 43 & sunlamp & 23.0 (21.5--24.6) & 0.1217 & 3.94 \\
    frozen & -- & 44 & in-domain test & 89.9 (87.9--91.7) & 0.1176 & 1.01 \\
    frozen & -- & 44 & synthetic (val) & 89.4 (88.0--90.7) & 0.1176 & 1.03 \\
    frozen & -- & 44 & lightbox & 50.6 (49.4--51.8) & 0.1176 & 2.50 \\
    frozen & -- & 44 & sunlamp & 22.1 (20.5--23.7) & 0.1176 & 3.74 \\
    \midrule
    frozen & \checkmark & 42 & in-domain test & 88.9 (86.8--90.8) & 0.1253 & 1.05 \\
    frozen & \checkmark & 42 & synthetic (val) & 89.6 (88.2--90.9) & 0.1253 & 1.01 \\
    frozen & \checkmark & 42 & lightbox & 66.2 (65.0--67.3) & 0.1253 & 1.78 \\
    frozen & \checkmark & 42 & sunlamp & 15.7 (14.3--17.1) & 0.1253 & 4.20 \\
    frozen & \checkmark & 43 & in-domain test & 86.2 (83.9--88.3) & 0.1259 & 1.10 \\
    frozen & \checkmark & 43 & synthetic (val) & 88.5 (87.0--89.9) & 0.1259 & 1.04 \\
    frozen & \checkmark & 43 & lightbox & 61.7 (60.5--62.9) & 0.1259 & 1.94 \\
    frozen & \checkmark & 43 & sunlamp & 13.4 (12.1--14.7) & 0.1259 & 4.46 \\
    frozen & \checkmark & 44 & in-domain test & 88.6 (86.5--90.5) & 0.1306 & 1.04 \\
    frozen & \checkmark & 44 & synthetic (val) & 89.5 (88.1--90.8) & 0.1306 & 1.02 \\
    frozen & \checkmark & 44 & lightbox & 67.8 (66.7--68.9) & 0.1306 & 1.69 \\
    frozen & \checkmark & 44 & sunlamp & 19.1 (17.7--20.6) & 0.1306 & 3.82 \\
    \midrule
    layer 4 & -- & 42 & in-domain test & 91.0 (89.1--92.7) & 0.0787 & 0.96 \\
    layer 4 & -- & 42 & synthetic (val) & 90.2 (88.8--91.5) & 0.0787 & 0.99 \\
    layer 4 & -- & 42 & lightbox & 34.1 (32.9--35.2) & 0.0787 & 6.07 \\
    layer 4 & -- & 42 & sunlamp & 13.5 (12.2--14.8) & 0.0787 & 6.93 \\
    layer 4 & -- & 43 & in-domain test & 92.0 (90.1--93.6) & 0.0739 & 0.93 \\
    layer 4 & -- & 43 & synthetic (val) & 91.3 (90.0--92.5) & 0.0739 & 0.95 \\
    layer 4 & -- & 43 & lightbox & 35.5 (34.4--36.7) & 0.0739 & 4.70 \\
    layer 4 & -- & 43 & sunlamp & 9.0 (8.0--10.2) & 0.0739 & 8.07 \\
    layer 4 & -- & 44 & in-domain test & 89.7 (87.6--91.5) & 0.0678 & 1.01 \\
    layer 4 & -- & 44 & synthetic (val) & 89.2 (87.8--90.6) & 0.0678 & 1.03 \\
    layer 4 & -- & 44 & lightbox & 32.2 (31.1--33.3) & 0.0678 & 5.11 \\
    layer 4 & -- & 44 & sunlamp & 12.0 (10.9--13.3) & 0.0678 & 7.68 \\
    \midrule
    layer 4 & \checkmark & 42 & in-domain test & 88.8 (86.7--90.7) & 0.0814 & 1.07 \\
    layer 4 & \checkmark & 42 & synthetic (val) & 90.5 (89.1--91.7) & 0.0814 & 0.99 \\
    layer 4 & \checkmark & 42 & lightbox & 43.7 (42.5--44.9) & 0.0814 & 4.58 \\
    layer 4 & \checkmark & 42 & sunlamp & 5.2 (4.4--6.1) & 0.0814 & 7.99 \\
    layer 4 & \checkmark & 43 & in-domain test & 90.3 (88.3--92.1) & 0.0911 & 1.00 \\
    layer 4 & \checkmark & 43 & synthetic (val) & 90.8 (89.5--92.1) & 0.0911 & 0.98 \\
    layer 4 & \checkmark & 43 & lightbox & 46.3 (45.1--47.5) & 0.0911 & 4.48 \\
    layer 4 & \checkmark & 43 & sunlamp & 5.7 (4.9--6.6) & 0.0911 & 6.96 \\
    layer 4 & \checkmark & 44 & in-domain test & 90.5 (88.5--92.2) & 0.0894 & 0.99 \\
    layer 4 & \checkmark & 44 & synthetic (val) & 91.0 (89.7--92.2) & 0.0894 & 0.97 \\
    layer 4 & \checkmark & 44 & lightbox & 51.4 (50.2--52.6) & 0.0894 & 3.60 \\
    layer 4 & \checkmark & 44 & sunlamp & 18.7 (17.2--20.2) & 0.0894 & 5.74 \\
    \bottomrule
  \end{tabular}
\end{table}

\begin{table}[ht]
  \caption{Rotation coverage (\%) at nominal 90\% for the ablation arms (per seed) and the frozen-backbone heads (range over seeds 42--44), each calibrated on its own predictions. $\hat{q}_R$ in rad.}
  \label{tab:conformal-revision}
  \centering
  \scriptsize
  \setlength{\tabcolsep}{4pt}
  \begin{tabular}{l c cccc c}
    \toprule
    Run & Seed & In-domain test & Synthetic (val) & Lightbox & Sunlamp & $\hat{q}_R$ \\
    \midrule
    A: no blur/noise & 42 & 90.5 & 91.1 & 33.2 & 16.0 & 0.558 \\
     & 43 & 89.2 & 88.8 & 33.9 & 16.8 & 0.538 \\
     & 44 & 90.4 & 90.5 & 31.8 & 18.7 & 0.540 \\
    \midrule
    B: jitter, gamma, erasing & 42 & 89.7 & 89.0 & 30.7 & 14.5 & 0.478 \\
     & 43 & 91.4 & 90.0 & 36.1 & 14.4 & 0.539 \\
     & 44 & 90.5 & 90.2 & 28.7 & 12.6 & 0.481 \\
    \midrule
    heads: 4k, logits & 42--44 & 89.1--92.4 & 89.5--92.2 & 85.3--87.1 & 83.0--85.3 & 2.81--2.86 \\
    heads: 4k, pooled & 42--44 & 88.5--91.9 & 88.5--92.5 & 77.0--83.2 & 71.0--77.7 & 2.57--2.71 \\
    heads: 28k, logits & 42--44 & 87.7--91.4 & 89.8--92.0 & 74.3--80.4 & 67.5--76.9 & 2.50--2.67 \\
    heads: 28k, pooled & 42--44 & 88.7--91.8 & 88.4--92.0 & 63.6--72.2 & 53.4--64.5 & 2.16--2.35 \\
    heads: tuned & 42--44 & 88.6--90.4 & 89.5--92.0 & 64.1--72.7 & 52.0--61.8 & 2.15--2.31 \\
    heads: synthetic BN & 42--44 & 89.4--90.1 & 89.3--90.5 & 73.6--77.1 & 72.2--74.2 & 2.64--2.67 \\
    \bottomrule
  \end{tabular}
\end{table}


\section{Constant-predictor baseline}
\label{app:constant}

The chance level of $\pi/2 + \nicefrac{2}{\pi} \approx 2.207$\,rad assumes orientations drawn
uniformly from $SO(3)$, which SPEED+ poses are not. This appendix gives the empirical
counterpart quoted in Section~\ref{sec:setup}: the score of a single fixed pose. It anchors
``near chance'' to a predictor one could actually build, rather than to an assumption about
the distribution.

The constant is fitted on the 28{,}000 training labels only. Fitting it on a test domain
would make it an oracle rather than a baseline, so neither real domain enters the estimate.
For rotation we take the principal eigenvector of $M = \sum_i \mathbf{q}_i
\mathbf{q}_i^{\top}$, which minimises chordal distance and is unaffected by the quaternion
double cover. We then refine it with five Karcher iterations on the sphere: sign-align each
$\mathbf{q}_i$ to the current mean, log-map into the tangent space, average, and exp-map back.
The $SO(3)$ geodesic distance is twice the $S^3$ one, so iterating on the sphere targets the
same minimiser as our reported metric. For translation we take the arithmetic mean of the
training translations. This gives $\mathbf{q} = (-0.198, 0.034, -0.019, -0.979)$ up to sign
and $\mathbf{t} = (-0.0001, 0.0026, 6.120)$\,m.

Two caveats belong with the numbers. First, the Karcher steps do not contract sharply here,
because the training orientations cover most of the sphere, which leaves the Fr\'echet
functional close to flat and the mean weakly identified. This does not affect the result: the
chordal and
Karcher estimates lie $0.022$\,rad apart and agree to $0.001$ on every score in
Table~\ref{tab:constant}. Second, both estimators minimise \emph{squared} error, whereas we
report mean unsquared error, whose minimisers are the corresponding geometric medians. The
scores below are therefore achievable rather than optimal, and they are an upper bound on the
error of the best constant predictor. We do not claim optimality, and no conclusion in the
paper needs it.

\begin{table}[ht]
  \caption{A constant pose fitted on the training split, evaluated on each test
  domain. $\Etrans$ range-normalised, $\Erot$ in radians. Skill is
  $2.207 - \Erot$. Reported for the Karcher-refined estimate.}
  \label{tab:constant}
  \centering
  \small
  \begin{tabular}{l c cc c}
    \toprule
    Domain & $n$ & $\Etrans$ & $\Erot$ & Skill \\
    \midrule
    synthetic (val) & 2{,}000 & 0.415 & 2.176 & 0.031 \\
    lightbox & 6{,}740 & 0.367 & 2.203 & 0.005 \\
    sunlamp & 2{,}791 & 0.372 & 2.199 & 0.009 \\
    \bottomrule
  \end{tabular}
\end{table}

The empirical no-skill rotation level therefore sits at most $0.031$\,rad from the uniform
one on all three domains, and we use $2.207$ in the body. Every trained run behind
Table~\ref{tab:results} keeps some skill over it, even the weakest, seed 42 of the frozen
un-augmented arm on sunlamp, which scores $2.048$\,rad and keeps $0.15$\,rad.

Translation is less comfortable, and it is where the second caveat matters most. On sunlamp
the constant scores $0.372$, which two of the three augmented layer-4 seeds ($0.405$, $0.399$)
do not beat, while a constant fitted directly to the unsquared metric does better on
translation, scoring
$0.330$ on sunlamp, and matches this one on rotation to within $0.001$. We quote the
Karcher-refined constant throughout, as the workshop version did, because no conclusion
changes with the tighter one: the same two seeds fail to beat it, and every other comparison
the paper draws with the constant still holds.


\section{Revision experiments}
\label{app:revision}

All runs in this appendix were added after review and pre-registered: predictions and
decision rules were written down and committed before any result was opened
(App.~\ref{app:prereg}), and selection uses synthetic validation only, as in the body.

\paragraph{Frozen-backbone block.} The frozen backbone runs in inference mode, so without
augmentation its pooled features (and the ImageNet logits computed from them) can be cached
once, and heads retrained on them in minutes. Applying the published seed-42 head to the
cache reproduces its published errors to within $0.0005$. Table~\ref{tab:frozen} gives three
seeds (42, 43, 44) for each configuration, all with the published head, loss and selection
rule. The 4{,}000-image cells use one fixed draw from the 28{,}000, not the original
4{,}000-image set. The tuned head is the best of learning rate $\{3{\times}10^{-4}, 10^{-3},
3{\times}10^{-3}\}\times\{20, 60\}$ epochs on seed 42, chosen on synthetic validation. The
synthetic-BN head re-estimates every BatchNorm layer on the un-augmented synthetic training
images, with the weights untouched, before caching.


\paragraph{Frozen backbone with augmentation.} Augmented inputs cannot use the feature cache,
so seeds 43 and 44 of the frozen augmented arm were trained end to end with the released code
on an Apple M5 GPU, while seed 42 is the published Colab run. Table~\ref{tab:frozenaug} pairs
each seed with the head-only retrain of the same seed in the published configuration
(Table~\ref{tab:frozen}), and their conformal coverage is in App.~\ref{app:conformal}.


\paragraph{Augmentation ablation.} Arm A is the full stack without blur and noise. Arm B keeps
only brightness and contrast jitter, gamma and random erasing (no glare blob, blur or noise).
Both train layer 4 on seeds 42, 43 and 44 with the published configuration, and we read them
as per-seed differences. They were trained with the released code on an Apple M5 GPU, whereas
the published arms were trained on Colab. A control rerun of the full stack (seed 42) on the
M5 differs from its published counterpart by $-0.053$, $+0.006$ and $-0.142$\,rad in $\Erot$
(synthetic, lightbox, sunlamp). That is within seed variation, but on sunlamp it is as large
as the ablation effects themselves.


\paragraph{Test-time BatchNorm adaptation (AdaBN).} For each layer-4 checkpoint we
re-estimate only layer 4's BatchNorm statistics, the ones fitted on synthetic data. We use
the unlabelled images of one half of each real domain (a fixed random split by filename),
leave every weight untouched, and evaluate on the other half (3{,}370 lightbox and 1{,}396
sunlamp images), with each arm's published synthetic-calibrated $\hat{q}$ for coverage. This is
test-time adaptation: it uses target images, though not their labels. The frozen arms are
excluded, as all their statistics are ImageNet's.


\paragraph{Image statistics.} Figure~\ref{fig:crops} compares pooled intensity densities,
which cannot tell a soft bright blob from a hard edge. Table~\ref{tab:imagestats} adds spatial
statistics, pooled over every image of each split after the same resize and grayscale
conversion (augmentation seed 7, as in Figure~\ref{fig:crops}):
\begin{itemize}[leftmargin=1.2em,itemsep=0pt,topsep=2pt]
  \item the Sobel gradient magnitude, scaled so a full black-to-white step reads 1
  \item edge density, the share of pixels whose gradient exceeds 0.1
  \item the clipped fraction, the share of pixels at or above $254/255$
  \item the total variation (TV) between pooled densities (48 intensity bins, 64 gradient
  bins)
\end{itemize}
The intensity TV of the full stack reproduces Figure~\ref{fig:crops} ($0.63$, $0.80$).


\section{Rotation-error shape}
\label{app:rotshape}

For each layer-4 checkpoint and real domain we measure two things, both pre-registered
(App.~\ref{app:prereg}). The first is the share of rotation errors within 0.3\,rad of $\pi$.
Under uniform guessing the angle density, $(1-\cos\theta)/\pi$, rises towards $\pi$, so chance
alone puts 0.19 of the mass there. The second is how concentrated the axes of those near-$\pi$
errors are, in the body frame ($R_\mathrm{true}^{\top} R_\mathrm{pred}$) and in the camera
frame ($R_\mathrm{pred} R_\mathrm{true}^{\top}$). We measure it as the share of axes within
$20^{\circ}$ of their principal axis, the leading eigenvector of the axes' second-moment
matrix. A flip about a near-symmetry axis would concentrate the body-frame axes, whereas an
in-plane confusion about the boresight would concentrate the camera-frame ones. Because the principal
axis is estimated from the same sample, the share's chance level is not the 0.06 of a fixed
axis. Simulation with uniform random axes gives 0.074--0.079 at these sample sizes (95th
percentile about 0.10).

\begin{table}[ht]
  \caption{Shape of the rotation errors on the real domains, layer-4 checkpoints. Near-$\pi$: share of errors within 0.3\,rad of $\pi$ (uniform guessing: 0.19). Cone: share of the near-$\pi$ errors' axes within $20^{\circ}$ of their principal axis, in the body and camera frames. A symmetry flip would concentrate the first and an in-plane confusion the second.}
  \label{tab:rotshape}
  \centering
  \scriptsize
  \setlength{\tabcolsep}{4pt}
  \begin{tabular}{l l c c c cc}
    \toprule
    Run & Domain & Median $\Erot$ & Near-$\pi$ & $n_{\pi}$ & Cone, body & Cone, camera \\
    \midrule
    layer 4, no aug., 42 & lightbox & 0.92 & 0.051 & 344 & 0.105 & 0.122 \\
     & sunlamp & 1.31 & 0.086 & 240 & 0.083 & 0.071 \\
    layer 4, no aug., 43 & lightbox & 0.92 & 0.050 & 339 & 0.198 & 0.112 \\
     & sunlamp & 1.54 & 0.095 & 265 & 0.132 & 0.068 \\
    layer 4, no aug., 44 & lightbox & 0.90 & 0.044 & 299 & 0.100 & 0.094 \\
     & sunlamp & 1.37 & 0.082 & 229 & 0.140 & 0.118 \\
    layer 4, full stack, 42 & lightbox & 0.74 & 0.027 & 183 & 0.191 & 0.131 \\
     & sunlamp & 1.50 & 0.090 & 250 & 0.128 & 0.092 \\
    layer 4, full stack, 43 & lightbox & 0.71 & 0.040 & 270 & 0.163 & 0.144 \\
     & sunlamp & 1.11 & 0.071 & 198 & 0.192 & 0.096 \\
    layer 4, full stack, 44 & lightbox & 0.69 & 0.039 & 263 & 0.171 & 0.122 \\
     & sunlamp & 1.31 & 0.071 & 197 & 0.315 & 0.071 \\
    \bottomrule
  \end{tabular}
\end{table}


\section{Coverage by range bin}
\label{app:range}

For each run we calibrate $\hat{q}$ on the calibration split exactly as in
Sec.~\ref{sec:conformal}, hold it fixed, and measure coverage within five range bins. The bin
edges are quintiles of the calibration split's range, so no test image shapes them. In-domain
bins hold 173--234 images each, and the per-bin counts sum to the marginal coverage on every
run. Tables~\ref{tab:rangebins-translation-metres}, \ref{tab:rangebins-translation-relative}
and~\ref{tab:rangebins-rotation} give each arm's coverage per bin, averaged over its seeds, for
the metre sphere, the range-normalised sphere and the attitude cone.

\begin{table}[ht]
  \caption{Coverage (\%) of the marginal 90\% region by range bin (m), translation, metres, as the mean over the seeds of each arm. Bins are quintiles of the calibration split's range.}
  \label{tab:rangebins-translation-metres}
  \centering
  \scriptsize
  \setlength{\tabcolsep}{4pt}
  \begin{tabular}{l l ccccc}
    \toprule
    Arm & Domain & $<3.8$ & 3.8--5.3 & 5.3--6.7 & 6.7--8.3 & $\geq8.3$ \\
    \midrule
    frozen, no aug. & in-domain & 97.5 & 95.7 & 90.2 & 82.9 & 75.2 \\
     & lightbox & 73.3 & 70.7 & 50.5 & 37.4 & 28.6 \\
     & sunlamp & 63.6 & 53.0 & 22.2 & 6.6 & 1.4 \\
    \midrule
    frozen, full stack & in-domain & 97.4 & 95.0 & 92.1 & 90.9 & 75.6 \\
     & lightbox & 91.3 & 84.1 & 64.4 & 58.2 & 45.6 \\
     & sunlamp & 81.7 & 39.4 & 11.6 & 3.0 & 0.1 \\
    \midrule
    layer 4, no aug. & in-domain & 92.9 & 95.0 & 92.7 & 89.8 & 80.9 \\
     & lightbox & 58.2 & 39.1 & 27.8 & 22.8 & 24.9 \\
     & sunlamp & 58.4 & 23.1 & 3.5 & 0.6 & 0.1 \\
    \midrule
    layer 4, full stack & in-domain & 92.0 & 94.1 & 93.6 & 89.3 & 84.8 \\
     & lightbox & 81.0 & 53.6 & 39.2 & 33.4 & 33.9 \\
     & sunlamp & 57.7 & 16.0 & 3.4 & 0.8 & 0.0 \\
    \bottomrule
  \end{tabular}
\end{table}

\begin{table}[ht]
  \caption{Coverage (\%) of the marginal 90\% region by range bin (m), translation, range-normalised, as the mean over the seeds of each arm. Bins are quintiles of the calibration split's range.}
  \label{tab:rangebins-translation-relative}
  \centering
  \scriptsize
  \setlength{\tabcolsep}{4pt}
  \begin{tabular}{l l ccccc}
    \toprule
    Arm & Domain & $<3.8$ & 3.8--5.3 & 5.3--6.7 & 6.7--8.3 & $\geq8.3$ \\
    \midrule
    frozen, no aug. & in-domain & 77.3 & 93.5 & 89.0 & 90.7 & 95.7 \\
     & lightbox & 38.7 & 56.2 & 50.1 & 48.5 & 54.0 \\
     & sunlamp & 31.8 & 39.6 & 21.9 & 9.7 & 5.3 \\
    \midrule
    frozen, full stack & in-domain & 74.5 & 88.5 & 88.6 & 94.3 & 92.9 \\
     & lightbox & 66.3 & 69.6 & 60.1 & 66.6 & 63.1 \\
     & sunlamp & 45.0 & 22.2 & 9.8 & 4.0 & 0.6 \\
    \midrule
    layer 4, no aug. & in-domain & 73.9 & 91.9 & 95.8 & 96.4 & 96.4 \\
     & lightbox & 35.1 & 33.7 & 31.1 & 32.3 & 38.3 \\
     & sunlamp & 35.0 & 18.0 & 4.6 & 1.0 & 0.2 \\
    \midrule
    layer 4, full stack & in-domain & 67.7 & 89.5 & 96.0 & 97.4 & 98.3 \\
     & lightbox & 57.0 & 48.3 & 42.8 & 41.9 & 47.0 \\
     & sunlamp & 33.7 & 12.3 & 3.6 & 1.0 & 0.1 \\
    \bottomrule
  \end{tabular}
\end{table}

\begin{table}[ht]
  \caption{Coverage (\%) of the marginal 90\% region by range bin (m), rotation, as the mean over the seeds of each arm. Bins are quintiles of the calibration split's range.}
  \label{tab:rangebins-rotation}
  \centering
  \scriptsize
  \setlength{\tabcolsep}{4pt}
  \begin{tabular}{l l ccccc}
    \toprule
    Arm & Domain & $<3.8$ & 3.8--5.3 & 5.3--6.7 & 6.7--8.3 & $\geq8.3$ \\
    \midrule
    frozen, no aug. & in-domain & 96.1 & 91.4 & 93.6 & 87.3 & 83.3 \\
     & lightbox & 75.9 & 74.2 & 71.1 & 69.3 & 64.2 \\
     & sunlamp & 64.2 & 56.5 & 61.4 & 56.2 & 52.5 \\
    \midrule
    frozen, full stack & in-domain & 92.0 & 92.8 & 92.9 & 90.1 & 84.9 \\
     & lightbox & 78.9 & 78.2 & 76.0 & 73.5 & 71.4 \\
     & sunlamp & 74.2 & 70.3 & 70.0 & 70.0 & 67.1 \\
    \midrule
    layer 4, no aug. & in-domain & 86.0 & 93.9 & 95.0 & 91.7 & 81.8 \\
     & lightbox & 22.8 & 29.0 & 27.5 & 24.5 & 19.7 \\
     & sunlamp & 11.1 & 11.6 & 11.3 & 10.3 & 7.1 \\
    \midrule
    layer 4, full stack & in-domain & 86.2 & 95.5 & 95.4 & 92.8 & 87.7 \\
     & lightbox & 35.6 & 40.9 & 38.3 & 34.5 & 29.5 \\
     & sunlamp & 17.3 & 17.7 & 16.8 & 11.0 & 9.3 \\
    \bottomrule
  \end{tabular}
\end{table}


\section{Reproducibility}
\label{app:repro}

\paragraph{Code, results and predictions.} Everything is in the public repository,
\url{https://github.com/lokindoo/spacecraft-pose-estimation}. The tag \texttt{arxiv-v1}
reproduces this version, and the tag \texttt{sim2science-2026-camera-ready} the workshop
version. Zenodo archives every tag under
\href{https://doi.org/10.5281/zenodo.24011140}{doi:10.5281/zenodo.24011140}, which lists each
version, and the workshop version alone under
\href{https://doi.org/10.5281/zenodo.24011141}{doi:10.5281/zenodo.24011141}. Each tag holds its
paper's source, with a script that checks every number in the paper against the repository's
\texttt{results/}. Per-image predictions for all 20 runs
behind the tables are released as CSV, so a new metric needs neither the data pipeline nor a
GPU.

\paragraph{Runs.} Configurations, seeds and the selection rule are as in
Sec.~\ref{sec:setup}. The eight published runs (seed 42 of both frozen arms, and all six
layer-4 runs) were trained on a free-tier Colab GPU. Frozen augmented seeds 43 and 44,
ablation arms A and B (three seeds each) and the control were trained with the released
package on an Apple M5 GPU. The frozen-backbone heads were trained on cached features on the
same machine's CPU, in minutes each.
On Colab a 20-epoch run took about 1.8\,h for the frozen un-augmented arm and 2.0--2.3\,h for
the layer-4 un-augmented arm. With augmentation, epochs took about 790\,s, some 4.4\,h per
run, whether or not the backbone was frozen. On the M5 most runs took 3.1--4.2\,h, and two
arm-A runs took about 14\,h of wall-clock time. A head on cached features trains in about
2.5 minutes.

\paragraph{Checkpoints.} \texttt{results/checkpoints.json} records, for each of the 17 trained
checkpoints, its configuration, seed, selected epoch, hardware and sha256. The weights are on
Zenodo with a model card, under CC~BY~4.0:
\href{https://doi.org/10.5281/zenodo.24140385}{doi:10.5281/zenodo.24140385}.

\paragraph{Software.} \texttt{requirements.txt} pins the versions the pipeline was written
against (PyTorch 2.13.0, torchvision 0.28.0). The M5 runs used Python 3.14.8, PyTorch 2.14.1
and torchvision 0.29.1, frozen in \texttt{requirements-m5.txt}. The Colab runs used the Colab
image of July 2026, whose versions were not recorded. The control rerun
(Sec.~\ref{sec:directional}) bounds what hardware and software together moved seed 42, and
none of the 81 unit tests needs data.

\section{Pre-registration record}
\label{app:prereg}

Every experiment added after review had its prediction and decision rule committed to version
control before its result was opened. The two files are published verbatim in the
repository's \texttt{preregistration/}, with every later edit and its time. The revision
predictions were committed on 28 Oct 2026 at 17:42 UTC. The extension's decision rules
followed on 30 Oct at 11:31 and its predictions at 11:54, and its results were opened at 13:36
the same day. These are commit times in the author's own repository, not a third-party
registration.

\begin{table}[ht]
  \caption{Pre-registered decision rules and their outcomes. Differences are paired by
  seed and given in rad, and ``3/3'' means all three seeds share the sign.}
  \label{tab:prereg}
  \centering
  \scriptsize
  \setlength{\tabcolsep}{3pt}
  \begin{tabular}{p{0.2\linewidth} p{0.38\linewidth} p{0.36\linewidth}}
    \toprule
    Experiment & Rule & Outcome \\
    \midrule
    \multicolumn{3}{@{}l}{\emph{Revision (locked 28 Oct)}} \\
    Arm A & Sunlamp, A $-$ full stack, 3/3 negative: keep ``blur and noise work against sunlamp outright'' & 2/3 ($-0.184$, $+0.006$, $-0.117$): softened to a hypothesis \\
    Arm B & Sunlamp, B $-$ A, 3/3 negative: the glare sentence gains evidence & 1/3 ($-0.018$, $+0.077$, $+0.072$): softened \\
    Control & Within 2 s.d.\ of published seed 42 on both real domains & Lightbox $+0.006$ ($\pm0.020$), sunlamp $-0.142$ ($\pm0.239$): inside; pairing holds \\
    Frozen seeds & Sunlamp s.d.\ $\geq 0.025$: ``only 0.03\,rad'' becomes ``no detectable change'' & s.d.\ 0.031: reworded \\
    Decomposition & One factor $\geq 2/3$ of the synthetic gain: name it; either factor moves sunlamp by $>2$ frozen s.d.: split the claim & Data 0.33, features 0.29 (53\%); largest sunlamp move 0.024 against 0.062: neither \\
    Head sweep & Beats published seed 42 by $>2$ frozen s.d.: ``under-trained''; recovers $\geq 50\%$ of the layer-4 gain on a real domain: restate & 0.913 against 1.053: tuned head reported; real domains unmoved: finding stands \\
    Synthetic-BN head & Share of the layer-4 gain recovered: $\geq 50\%$ / 0--50\% / $\leq 0$ & $-0.51$ (lightbox), $-0.27$ (sunlamp): ``not explained by normalisation statistics alone'' \\
    AdaBN & Improved if 6/6 checkpoints share the sign; coverage $<80\%$: ``coverage stays low'' & $\Erot$ 6/6 on both domains; $\Etrans$ 5/6 lightbox, 6/6 sunlamp; coverage at most 40.4\% \\
    Image statistics & Sunlamp clipped more than the full stack, and sharper than lightbox and the full stack; else two passages become hypotheses & Clipping holds (2.78\% vs 1.96\%); edge density fails (2.0\% vs 6.6\%, 10.7\%): passages reworded \\
    \midrule
    \multicolumn{3}{@{}l}{\emph{Extension (locked 30 Oct)}} \\
    Frozen + aug & Synthetic aug $-$ no aug, 3/3 positive: keep ``the reference degrades''; sunlamp 3/3 negative and larger than lightbox: qualify the directional claim & $+0.169$, $+0.161$, $+0.134$: kept; sunlamp $-0.128$, $-0.069$, $-0.132$, larger than lightbox on every seed: qualified \\
    Error shape & Near-$\pi$ share $\geq 0.28$ on all six: bimodal; body cone $\geq 0.18$ on all six and camera not: flips; reverse: in-plane & Near-$\pi$ 0.03--0.10: not bimodal; body cone $\geq 0.18$ on 2/6 per domain: neither \\
    Range bins & Both in-domain translation patterns hold and the extreme bin excludes 90\%: keep; monotone but intervals include 90\%: soften; not monotone: cut & Ambiguous: ``monotone'' undefined. Near-to-far, both hold and the extreme bin excludes 90\% on 7/8 and 8/8; strictly over five bins, 3/8 and 2/8. Reported as near and far numbers \\
    Ablation coverage & Coverage opposite to error on $\geq 5/6$ per domain: ``coverage mirrors error'' & 4/6 on each domain: not claimed \\
    \bottomrule
  \end{tabular}
\end{table}

\paragraph{Deviations and corrections.} Frozen augmented seeds 43 and 44 were pre-registered
for this version but finished after the camera-ready's results freeze. Once their outcome was
known they were added to the camera-ready anyway, the only result to cross the freeze. The
range-bin rule also left ``monotone'' undefined, as Table~\ref{tab:prereg} notes.

\paragraph{Not pre-registered.} Three results in the body were found or computed outside the
pre-registration: the transfer plot (Fig.~\ref{fig:transfer}), AdaBN's effect on translation (Sec.~\ref{sec:bn}) and the collapse
of sunlamp translation coverage with range (Sec.~\ref{sec:conformal}). Every point in the
transfer plot had already been published.
